\documentclass[11pt,a4paper]{article}

\usepackage[margin=1.8cm]{geometry}
\usepackage[T1]{fontenc}
\usepackage[utf8]{inputenc}
\usepackage{lmodern}
\usepackage[hidelinks]{hyperref}
\usepackage{enumitem}
\usepackage{titlesec}
\usepackage{tabularx}
\usepackage{xcolor}
\usepackage{array}

\pagestyle{empty}
\setlength{\parindent}{0pt}
\setlength{\parskip}{0.25em}
\setlist[itemize]{leftmargin=1.2em,itemsep=0.2em,topsep=0.25em}

\titleformat{\section}{\large\bfseries}{}{0em}{}[\vspace{-0.35em}\titlerule]
\titlespacing*{\section}{0pt}{1.0em}{0.45em}

\newcommand{\nameheader}[3]{%
  {\LARGE\bfseries #1}\\[0.35em]
  #2 \enspace $\cdot$ \enspace \href{mailto:#3}{#3}\\[0.2em]
}

\newcommand{\cventry}[4]{%
  \begin{tabularx}{\textwidth}{@{}Xr@{}}
    {\textbf{#1}} & {#4}\\
    {\emph{#2}} & {#3}
  \end{tabularx}\vspace{0.2em}
}

\newcommand{\cvsubentry}[2]{%
  \begin{tabularx}{\textwidth}{@{}Xr@{}}
    {\textit{#1}} & {#2}
  \end{tabularx}\vspace{0.15em}
}

\begin{document}

\nameheader{Simone Ghezzi Colombo}{Italy}{simoneghezzi24@gmail.com}

\section*{Education}

\cventry{Incoming PhD student, Analysis of Social and Economic Processes (ASEP)}{University of Milano-Bicocca}{Milan, Italy}{Starts Nov 2026}
Research interests: social and economic processes, inequality, local welfare, and quantitative methods.

\cventry{M.A. in Programming and Management of Social Policies and Services}{University of Milano-Bicocca}{Milan, Italy}{Oct 2023 -- Oct 2025}
Final grade: 110/110 \emph{cum laude}; GPA: 30.09/30. Thesis: \emph{Masculinities, Feminist Mobilization and the New Cultural Divide: A Quantitative Analysis of Gendered Political Attitudes Among Young Europeans.}
\begin{itemize}
  \item Built an original EU27 comparative dataset by harmonizing ESS Round 11 (2023--2024) and EES 2024 microdata, linking vote choice to CHES 2024 expert-coded party positions through a manually validated party-ID crosswalk.
  \item Designed a reproducible R workflow for data preparation, cohort trends, country comparisons, and multivariate modelling, with documented harmonization procedures, diagnostics, and version control.
  \item Produced figures and EU maps using centred rolling windows, LOESS smoothing, confidence intervals, and transparent trimming rules to reduce instability and outlier influence.
\end{itemize}

\textbf{Selected coursework research}
\begin{itemize}
  \item \textbf{Planning and Project Management:} designed an evidence-based intervention on youth-related challenges in Brianza using Project Cycle Management, Stata, and QGIS.
  \item \textbf{Globalization and Local Development:} conducted a comparative analysis of post-socialist EU member states using R, Python, Eurostat, ARDECO, and QGIS to map regional disparities.
  \item \textbf{Quantitative Analysis of Social Phenomena:} used Stata to analyse ISSP 2019 \emph{Social Inequality V}, studying the relationship between subjective class position and future expectations.
  \item \textbf{Welfare and Public Action:} examined the political dimensions of local welfare governance, with a focus on technicization, depoliticization, and democratic control.
  \item \textbf{Processes of European Integration:} explored the relationship between English proficiency, social media use, and European identity among young Europeans.
  \item \textbf{Social Change, New Social Risks and Social Innovation in Europe:} conducted a survey experiment on welfare chauvinism among students at the University of Milano-Bicocca.
  \item \textbf{Social Planning:} evaluated Free-Fare Public Transport policies through comparative case studies and literature review.
  \item \textbf{Systems of Family Solidarity:} analysed the geography of gender-based violence through cross-cultural and cross-national comparison.
\end{itemize}

\cventry{Erasmus Exchange Semester}{Universit\"at Konstanz}{Konstanz, Germany}{Sep 2020 -- Mar 2021}
\begin{itemize}
  \item \textbf{The Social Neuroscience of Empathy:} analysed altruistic punishment and empathic anger in contexts of unfair treatment.
  \item \textbf{European Integration Theories:} examined the 2020 shift in Italian politics from populist Euroscepticism to pro-European technocracy through a post-functionalist lens.
\end{itemize}

\cventry{B.A. in Sociology}{University of Milano-Bicocca}{Milan, Italy}{Sep 2019 -- Oct 2022}
Final grade: 110/110 \emph{cum laude}. Thesis: \emph{Queer Migration: The Role of Sexuality in Moving to a Foreign Country.}
\begin{itemize}
  \item Conducted research on studentification in European cities through a comparative analysis of Loughborough and Urbino.
  \item Investigated the role of social media in youth socialization through interviews and qualitative analysis.
  \item Developed a qualitative research project on queer migration, focusing on sexuality, mobility decisions, and post-relocation experiences.
\end{itemize}

\section*{Professional Experience}

\cventry{Open Data Observatory -- internship}{Department of Sociology and Social Research, University of Milano-Bicocca}{Milan, Italy}{Apr 2025 -- Jul 2025}
\begin{itemize}
  \item Contributed to the structuring, validation, and documentation of institutional datasets.
  \item Developed reproducible data pipelines in R and Git, ensuring methodological rigor, version control, and compliance with open-data standards.
  \item Assessed data quality, completeness, and interoperability across institutional sources to support policymaking and research.
\end{itemize}

\cventry{President of the Youth Council}{Municipality of Olgiate Molgora}{Olgiate Molgora, Italy}{Jul 2022 -- Present}
\begin{itemize}
  \item Coordinated educational street-art projects and developed a proposal to turn a local park into a social space for young people.
  \item Organized public initiatives and opportunities for community participation.
\end{itemize}

\cventry{Co-founder and Organizer}{EUatSchool}{Merate, Italy}{Feb 2024 -- Present}
\begin{itemize}
  \item Co-founded a non-partisan, student-led European civic education initiative. Its third edition engaged more than 300 upper-secondary students and over 20 young volunteers in direct exchanges with Italian and European institutions.
  \item Helped secure the patronage of the European Parliament and the European Commission Representation in Italy, together with support from local institutions, schools, and civil-society partners.
  \item Organized and moderated exchanges with MEPs and senators, and coordinated workshops and public debates on European integration, foreign policy, defence, energy, AI, and the digital euro.
\end{itemize}

\cventry{Project Coordinator}{EVOLVO SRL}{Milan, Italy}{Dec 2022 -- Jun 2024}
\begin{itemize}
  \item Managed communication with European partners, colleagues, and Erasmus+ National Agencies.
  \item Coordinated more than 100 Erasmus+ projects annually through structured information management and documentation.
  \item Organized multi-day team-building events for international groups and delivered orientations to more than 1,000 participants per year.
\end{itemize}

\cventry{Civil Servant}{Italian Red Cross}{Olgiate Molgora, Italy}{Apr 2021 -- Apr 2022}
\begin{itemize}
  \item Coordinated dialogue with municipalities, hospitals, and health institutions to support local service delivery and operational planning.
  \item Helped manage health and social services and coordinate more than 200 active volunteers.
\end{itemize}

\section*{Papers, Talks, and Academic Presentations}

\cventry{Paper accepted for presentation at the ECPR General Conference 2026}{European Consortium for Political Research (ECPR)}{Krak\'ow, Poland}{Sep 2026}
Single-authored paper accepted in the panel \emph{Gender, Democracy, and the Menace of Authoritarianism}. Paper title: \emph{From Gap to Cleavage: Gender Polarization Among Gen Z in the EU.}

\cventry{Invited guest lecture on youth political participation in Europe}{Department of Sociology and Social Research, University of Milano-Bicocca}{Milan, Italy}{Mar 2026}
Invited by Prof.~Tatjana Sekuli\'c to lecture in the Master's course \emph{Le Dinamiche delle Integrazioni Europee} (PROGEST, a.y.~2025/2026), presenting original research based on the Master's thesis.

\cventry{Speaker and co-organiser, public seminar on mobility and public transport}{Viviamo Merate}{Merate, Italy}{May 2026}
Presented an original analysis of the S8 suburban railway line, focusing on service evolution, passenger growth, and the role of rail in regional mobility planning.

\section*{Awards and Recognition}

\cventry{Selected as one of six thesis presenters, 3rd Convivio Progest 2026}{Department of Sociology and Social Research, University of Milano-Bicocca}{Milan, Italy}{Mar 2026}
Selected among six Master's graduates invited to present thesis research at a departmental event showcasing outstanding theses discussed in the previous 18 months.

\section*{Research Projects}

\cventry{Computational analysis of Italian parliamentary elections at municipal level}{Independent / GitHub-based research infrastructure}{Italy}{Mar 2026 -- Present}
\begin{itemize}
  \item Developing a reproducible and well-documented dataset for the study of electoral behaviour, territorial variation, and political competition.
  \item Implementing validation pipelines, derived indicators, and structured access tools for quantitative and computational social-science research.
  \item Building outputs that make political data explorable and reusable for future comparative work across countries, institutions, and political contexts.
\end{itemize}

\section*{Skills}

\textbf{Methods and tools:} R, Python, Stata, QGIS, Git/GitHub, data harmonization, reproducible workflows, quantitative analysis, survey research, comparative research, spatial analysis.\\
\textbf{Languages:} Italian (native), English (C2), Spanish (B1).

\section*{Academic and Civic Service}

\cventry{Volunteer and community organizer}{Italian Red Cross and local civic initiatives}{Olgiate Molgora, Italy}{Apr 2021 -- Present}
\begin{itemize}
  \item Coordinate community-based educational and social initiatives, including food and school-supply drives, activities for older adults, and inclusive events for children and families.
  \item Contribute to project planning, institutional communication, sponsorship requests, funding applications, and relations with local stakeholders.
\end{itemize}

\end{document}
